\chapter{Introduction}
\label{ch:intro}

\section{Motivation}
Large language model (LLM) agents are increasingly deployed on tasks whose
relevant information far exceeds a single prompt: answering a question that
spans hundreds of pages of financial filings, recalling a fact stated five
conversational sessions ago, or chaining evidence across several documents. The
dominant response is to bolt on a \emph{memory}: a context window of a fixed
size, a retrieval-augmented generation (RAG) index with a fixed chunking and
scoring rule, or a vector store that treats every observation alike.

What these solutions share is that the memory is \textbf{fixed}. The decision of
\emph{what to store}, \emph{which concepts to track}, and \emph{how to score a
retrieval} is hand-designed once and frozen. Yet different tasks plainly demand
different memory structures. A multi-hop question needs entity relations; a
navigation task needs spatial and temporal transitions; an end-of-corpus
question, asked after the agent has ingested an entire corpus, needs
\emph{recency-independent} recall, because the relevant fact may have been seen
long ago. A memory rule that is optimal for one of these is often actively
harmful for another.

A single parameter makes the tension concrete. Consider how much a retriever
should favour \emph{recent} evidence. A conversational assistant answering
``what did I just ask you to change?'' wants recency weighted heavily --- the
answer is almost always in the last few turns. An agent answering ``what is the
governing-law clause?'' after ingesting a five-hundred-page contract wants recency
weighted at essentially \emph{zero} --- the relevant clause was read hundreds of
pages ago, and a recency-biased retriever will confidently surface the most
recent (wrong) clause instead. The same is true of how aggressively to store
(store everything, and retrieval drowns in noise; store too little, and the
answer was never kept) and of how much to trust embedding similarity versus exact
lexical overlap. No single setting of these knobs is good everywhere; the right
setting is a property of the task. In practice this is ``solved'' by an engineer
hand-tuning the memory for each deployment --- ad hoc, unmeasured, and rarely
revisited. The premise of this thesis is that those knobs can instead be written
down explicitly and \emph{learned} per task, and that doing so both improves
retrieval and makes the task-dependence measurable rather than anecdotal.

\section{Research question}
This thesis asks a single question:
\begin{quote}
\itshape Can an agent \textbf{learn how to construct its own memory} --- what to
store, which concepts to track as nodes, and how to score retrieval --- and is
the learned optimum \textbf{task-dependent}?
\end{quote}
We make the construction of memory itself a learnable object. A small parameter
vector $\thetavec$ governs storage, abstraction, and retrieval; we then
\emph{optimize $\thetavec$} --- not the agent's policy and not the LLM's
weights --- with black-box search, and we ask whether the recovered $\thetavec$
differs by task and whether adapting it improves end-task performance.

\section{Contributions}
The contribution is \textbf{not} a better reinforcement-learning policy or a new
language model. It is the claim and the demonstration that \emph{memory
construction is a learnable, task-adaptive object}. Concretely, this thesis:
\begin{itemize}[leftmargin=1.4em]
  \item \textbf{Formalizes memory construction as a parameter vector
  $\thetavec$} (storage probability, entity/temporal thresholds, and learnable
  retrieval weights \wgraph{}/\wembed{}/\wrec{}, up to ten dimensions in the most
  complete system, \Vfour{}) and optimizes it by Evolution Strategy and CMA-ES.
  \item \textbf{Shows task-dependence at three increasing levels of realism}:
  grid worlds (Stage~1), a twelve-system $\times$ four-environment benchmark with
  ablation, transfer, and sensitivity analysis (Stage~2), and corpus-cumulative
  question answering with real LLMs over six long-context benchmarks (Stage~3).
  \item \textbf{Subjects its own headline results to an adversarial self-audit}
  and reports the corrected numbers --- including the finding that a
  full-context ``dump-all'' baseline, once a truncation bug was fixed, is
  statistically tied with selective retrieval on accuracy, which reframes the
  contribution of learned memory at this scale as one of \emph{efficiency}
  rather than raw accuracy.
\end{itemize}

\section{The central finding}
Across the three stages the same claim holds: \emph{what an agent stores and how
it scores retrieval can be learned, and the learned optimum is task-dependent}.
In the grids the optimizer recovers visibly different $\thetavec$ per task and
lifts the hardest task from $2.5\%$ to $27.5\%$ success. On the benchmark suite a
ten-dimensional learnable memory reaches the top performance cluster, and a
$\thetavec$ tuned for one long-haystack task generalizes within that task family
but not across families. On six real long-context benchmarks, re-tuning
$\thetavec$ on a corpus-cumulative objective yields a recency-independent memory
whose end-of-corpus accuracy lift \textbf{survives held-out testing on two of the
three domain-coherent benchmarks} (FinanceBench and CUAD; QASPER is
non-significant); two further benchmarks serve as pre-registered
domain-incoherent controls, and NarrativeQA's objective is undefined by
construction.

A second methodological result emerges from comparing learned selective memory
against a fixed full-context baseline (``dump everything into the prompt''): once
measured correctly, the two are \emph{accuracy-competitive}. The argument for
learnable, task-adaptive memory at this corpus scale is therefore \textbf{cost
and scalability} --- it matches accuracy at roughly one-eighteenth of the token
spend and does not break when the corpus exceeds the context window --- rather
than a claim that selective retrieval is uniquely accurate.

\section{Thesis outline}
Chapter~\ref{ch:background} reviews long-context retrieval, agent-memory
systems, and black-box optimization, and positions this work. Chapter~\ref{ch:method}
formalizes $\thetavec$, the graph-memory architecture, the optimizers, the three
experimental settings, and the evaluation protocol. Chapter~\ref{ch:results}
reports the experiments stage by stage. Chapter~\ref{ch:discussion} interprets
the results, discusses threats to validity, and summarizes the adversarial
self-audit. Chapter~\ref{ch:conclusion} concludes and lists future work. Four
appendices give the reproducibility material, the adversarial self-audit, the
complete per-benchmark tables, and the data-preparation and judging details.
